% AI-produced by GPT-6 Astra (OpenAI, ChatGPT Pro mode). No priority claim.
\subsection{A functorial contraction of the binary and triple actions}

\begin{proposition}[Proved: the path cone is a functor]
\label{A:predimension-cone-functor}
For either category $A$ of
Proposition~\ref{A:predimension-quasi}, there is an endofunctor
$F:A\to A$ with natural transformations
\[
 \eta:\operatorname{id}_A\longrightarrow F,
 \qquad \gamma:\operatorname{const}_V\longrightarrow F.
\]
The map $\eta_X$ includes $X$, and $\gamma_X$ selects a new centre
$c_X$. For every $V$-arrow $x:V\to X$, the two vertices
$\eta_Xx$ and $\gamma_X$ are joined by a directed binary path of
three actual $W$-arrows in $F(X)$.
\end{proposition}
\begin{proof}
Work first with a concrete category of finite structures, with a
chosen small set of labels closed under the following finite tagged
constructions; passage to a skeleton transports the functor and its
natural transformations. Use a tagged disjoint copy for the old
vertices as well, so all the following labels are disjoint
functorially. For every raw vertex $x$ of $X$, adjoin
five tagged points $r_x,s_x,a_x,b_x,d_x$, together with one new
centre $c_X$, and the three hyperedges
\[
 \{x,r_x,a_x\},\qquad \{r_x,s_x,b_x\},\qquad
 \{s_x,c_X,d_x\}.
\]
Give $F(X)$ the block order used in
Proposition~\ref{A:predimension-binary-colimit}: old points first;
then $r_x<s_x$ blocks ordered by $x$; then $c_X$; then
$a_x<b_x<d_x$ blocks ordered by $x$. The old object is $d$-closed
in $F(X)$ by first adjoining the centre disjointly, then adjoining
the paths using Lemma~\ref{A:predimension-three-path}.
The centre is $d$-closed in the disjoint union and hence in $F(X)$.

For an arrow $f:X\to Y$, define $F(f)$ on old points by $f$, send
$c_X$ to $c_Y$, and send each tagged point indexed by $x$ to the
point with the same tag indexed by $f(x)$. This is an increasing
induced embedding, and the formulas give functoriality exactly.
To prove it $d$-closed, identify $X$ with its image in $Y$.
The subhypergraph $Y\cup_X F(X)$ of $F(Y)$ is a free amalgam over
the $d$-closed object $X$. Therefore
\[
 F(X)\leq_d Y\cup_X F(X).
\]
Adjoin the remaining paths, indexed by $y\in Y\setminus X$,
one at a time. Their two endpoints already exist; the path lemma
makes each extension $d$-closed. Thus $F(X)\leq_d F(Y)$.
The displayed inclusion and centre maps are natural by their
tagged definitions.

It remains to verify that the required paths are paths of actual
morphisms, rather than merely raw hyperedges. A loose path has
every point and every edge $d$-closed. This follows by the same
calculation as Lemma~\ref{A:predimension-loose-cycle}: a chosen
set of binary vertices induces a forest, and its minimum
predimension after selecting private points is
$2|B'|-e(B')\geq |B'|+1$ when $B'$ is nonempty. In particular,
a proper extension of a point has predimension at least three,
and a proper extension of an edge has predimension at least four.

Suppose now that $x:V\to X$ is an actual arrow, so $\{x\}\leq_d X$.
Let $Q_x$ be its loose path together with its two endpoints $x,c_X$.
The union $X\cup_{\{x\}}Q_x$ is a free amalgam over a $d$-closed
point. Consequently $Q_x$ is $d$-closed in that union. Adjoining
all the other paths by the path lemma proves $Q_x\leq_d F(X)$.
All its edges are therefore actual $W$-arrows. Their directions
are increasing by the chosen block order.

If $X$ belongs to the point-and-edge-closed subcategory, this last
argument applies to every raw $x$. Every new point and edge then
lies in one of the $d$-closed paths $Q_x$; the old atoms remain
closed by $X\leq_d F(X)$. If the empty object is included, its
image is just the one-point object. Thus $F$ preserves that
subcategory too.
\end{proof}

\begin{theorem}[Proved: contractibility of the binary and triple element categories]
\label{A:predimension-cone-contracts}
Let $A$ be either category of
Proposition~\ref{A:predimension-quasi}. The categories of elements
of each of the two actions
\[
 H_{12}=A(V,-)/(ev_1\sim ev_2),
 \qquad
 H_{123}=A(V,-)/(ev_1\sim ev_2\sim ev_3)
\]
have contractible classifying spaces. More generally this holds for
every pointwise quotient $H$ of $H_{12}$. In particular all their
positive-degree integral homology groups vanish. For the binary
cycle functor $K$ of Proposition~\ref{A:predimension-binary-colimit},
\[
 T_i^A(K)=0\qquad(i>0).
\]
\end{theorem}
\begin{proof}
Use $F,\eta,\gamma$ from
Proposition~\ref{A:predimension-cone-functor}. The natural
transformations
$\operatorname{id}_A\to F\leftarrow\operatorname{const}_V$
make the identity of $BA$ homotopic to a constant map. The category
is nonempty, so $BA$ is contractible.

Let $H_{12}\twoheadrightarrow H$ be the given quotient. Write
$h_X\in H(FX)$ for the image of the centre arrow $\gamma_X$.
These elements are natural in $X$. Every element $h\in H(X)$ is
represented by an actual point arrow $x:V\to X$. The length-three
path from $\eta_Xx$ to $\gamma_X$ consists of actual $W$-arrows,
so its binary relations give
\[
 H(\eta_X)(h)=h_X.
 \tag{CC}\label{eq:predimension-central-collapse}
\]
This equality also holds in any further quotient, including
$H_{123}$.

Let $p:\operatorname{El}_A H\to A$ be the projection, and define
\[
 J:A\longrightarrow\operatorname{El}_A H,
 \qquad J(X)=(F(X),h_X),\qquad J(f)=F(f).
\]
Naturality of the centre makes $J$ a functor.
Equation~\eqref{eq:predimension-central-collapse} says that
$\eta_X$ is an arrow
\[
 (X,h)\longrightarrow (F(X),h_X)=Jp(X,h)
\]
in the category of elements. These arrows form a natural
transformation $\operatorname{id}_{\operatorname{El}H}\to Jp$.
After taking nerves, the identity therefore factors up to homotopy
through the contractible space $BA$. The element category is
nonempty, for example because $H(FV)$ contains its centre class.
Its identity is homotopic to a constant map, proving the asserted
contractibility.

Finally the binary presentation by two projective representables,
split at the image of its middle map, gives dimension shifting
\[
 T_i^A(K)\cong T_{i+2}^A(\Z[H_{12}])
       \cong H_{i+2}(B\operatorname{El}_A H_{12};\Z)=0
 \qquad(i>0).
\]
\end{proof}

\paragraph{An equivalent homological factorization.}
The cone also supplies a natural map
\[
 \Z[A(W,-)]\longrightarrow F^*K,
 \qquad e\longmapsto e+P_{ev_2}-P_{ev_1}.
\]
Its restriction along $K\hookrightarrow\Z[A(W,-)]$ is exactly
$K\xrightarrow{K\eta}F^*K$, because the path coefficients cancel
on a cycle. The standard derived-colimit comparison
$T_i(F^*K)\to T_i(K)$, induced by the indexing functor $F$,
composes with $T_i(K\eta)$ to the identity: the natural
transformation $\eta$ gives the usual bar-chain homotopy.
For $i>0$ the factorization through a projective representable
therefore gives the same vanishing. The direct nerve contraction
above additionally proves the homotopy statement.

\paragraph{Scope of the cone argument.}
The contraction uses an action generated by point arrows, together
with the binary hyperedge relation. An arbitrary $A$-set need not
be such a quotient. For a general quotient of a representable
$A(P,-)$, common fixed substructures of its presenting relations
can retain information which the point cone does not remove.
The passage to every action is supplied by the separate
finite-support argument culminating in
Theorem~\ref{A:predimension-total}.
